\subsection{Column Vectors as Matrices}

A column vector may be viewed as a matrix with one column. This allows vector and
matrix equations to be written in one language.

\begin{examplebox}[One vector, two orientations]
For
\[
u=
\begin{pmatrix}
2\\-1\\4
\end{pmatrix},
\]
we have a $3\times1$ matrix. Its transpose is the $1\times3$ row matrix
\[
u^{\mathsf T}=\begin{pmatrix}2&-1&4\end{pmatrix}.
\]
\end{examplebox}

\begin{interpretationbox}[Orientation changes what products mean]
The distinction between a row and a column is structural. The products
\[
u^{\mathsf T}u
\qquad\text{and}\qquad
uu^{\mathsf T}
\]
have different sizes: the first is $1\times1$, while the second is $3\times3$.
Even when the same entries occur, their arrangement determines which operations
are possible and what object is produced.
\end{interpretationbox}

\begin{remarkbox}[Read dimensions before multiplying]
The notation may look compact, but dimensions carry the compatibility check.
Always determine the shapes before applying the row-by-column rule.
\end{remarkbox}
